\section{Products and intrinsic memory growth}
\label{sec:tensor29}
A classification becomes more useful when it survives independent
composition. For the following statement a continuation is a normalized
channel from a finite future-query alphabet $Y$ to a finite output
alphabet $O$. Multiple components are allowed. The joint experiment
requires the whole product output law, not merely the correct marginal
means.

\begin{theorem}[Multiplicativity on the complete class]
\label{thm:tensor29}
Let $\mathscr R\subset\prod_{y\in Y}\Delta(O)$ and
$\mathscr S\subset\prod_{z\in Z}\Delta(A)$ be finite sets whose
support-component hulls are simplices. Form the independent product
rows
\[
 (R\otimes S)(o,a\mid y,z)=R(o\mid y)S(a\mid z).
\]
Their component hulls are again simplices, and their intrinsic capacities
satisfy
\begin{equation}\label{eq:tensor29}
                     s(\mathscr R\otimes\mathscr S)
                           =s(\mathscr R)s(\mathscr S).
\end{equation}
Consequently every joint stochastic encoder and decoder for the product
requires this many states; correlated internal implementations cannot
reduce the number. For independent terminal channels whose residual
sets pass the test at every individual cut, the product profile on
any interleaving is exactly the product of the individual capacities
at the corresponding vector of cuts.
\end{theorem}
\begin{proof}
Minimal faces of a query channel are given by the allowed outcomes at
each query. Two product faces intersect exactly when the faces in each
factor intersect: supports multiply, and a nonempty intersection must
have nonempty projected supports in both factors for every query.
The component graph of the product is the graph product with a loop
at each vertex. Its connected components are precisely products of
components, because one can move along a path in one factor while
holding the other fixed.

In a component with $k$ simplex vertices, the flattened normalized
channel vectors of those vertices are linearly independent. Indeed,
normalization turns any linear dependence into an affine dependence.
For components with $k$ and $l$ vertices, their $kl$ product vectors
are linearly independent by the tensor-product rank identity. Every
product of rows is a convex combination of these vectors, and every
one of these product vertices is itself a product row. Hence the
component hull is their $(kl-1)$-simplex. Summing $kl$ over pairs of
components gives \eqref{eq:tensor29}. The cut lower bound allows
arbitrary normalized joint generators, not just product generators.
The upper construction uses products of the individual generators.
Restriction of a product residual changes only the factor whose input
is read, so compatibility holds simultaneously throughout the
interleaving graph.
\end{proof}

The independence in the specification of the output channel is
indispensable. Independent inputs or additive payoffs alone do not
force a product of optimal response channels. In particular,
\eqref{eq:tensor29} is not a multiplicativity claim for arbitrary
nonnegative rank or for the marked minimax values $U_N$.

\begin{corollary}[A full-support family at a preparation barrier]
\label{cor:barrier29}
For $q\ge2$ and $0<\lambda<1$, let
\[
 T_\lambda(b\mid a)=\lambda\1_{a=b}+(1-\lambda)/q,
 \qquad a,b\in[q].
\]
Acquire $n$ input symbols before releasing any output; subsequently
produce the specified independent product channel $T_\lambda^{\otimes n}$.
The exact minimum peak is $q^n$, even if acquisition and output
schedules depend on retained observations and internal randomization,
provided all schedule memory is charged and no past transcript can be
read back. For a declared schedule, at quiescent cuts between completed release
operations, the exact width with $m$ acquired but unreleased symbols
is $q^m$. An intermediate held-output cut is charged separately using
the active count immediately before that release.
At $\lambda=0$, one state suffices until an output-buffer cut.
\end{corollary}
\begin{proof}
$T_\lambda$ has eigenvalues $1$ and $\lambda$ with multiplicity $q-1$,
so its $q$ rows are independent and form a simplex. All entries are
positive. Theorem~\ref{thm:tensor29} gives $q^m$ at any fixed active cut.
For an adaptive schedule at the barrier, condition on the complete
input word $x$. Its state law is an encoder $E(s\mid x)$; the law of
the entire subsequently produced labelled output vector is a decoder
$D(\mathbf b\mid s)$. Any output-order choices are included in this
law before taking the labelled-output projection. Thus
$T_\lambda^{\otimes n}=ED$, whose rank $q^n$ bounds the state count.
This argument charges, rather than reveals, the acquisition transcript.
Retaining the full input word gives the matching fixed-order upper
bound. Releasing a symbol samples its output and discards that input;
even if the output itself is held for one cut, there are at most
$q\,q^{m-1}=q^m$ states. At $\lambda=0$ the outputs are independent
of all inputs, and can be sampled without retaining them.
\end{proof}

This is a rank-boundary phenomenon entirely inside full support, unlike
the physical support boundary in Theorem~\ref{thm:reverse-jump}.
For every fixed $n$, the product law tends uniformly in total variation
to the input-independent output law, with defect at most
$n\lambda(1-1/q)$. Thus a one-state generator approximates it with that
defect, while every nonzero $\lambda$ has the stated exact memory.
The estimate follows by coupling each output with an independent
uniform symbol and summing the per-coordinate disagreements. Exact
memory and approximate statistical performance are different quantities.
